\documentstyle [11pt]{article}
\textheight 200mm
\title{Implicitization by Gr\"obner basis conversion}
\author{Michael Kalkbrener\\Department of Mathematics\\
Swiss Federal Institute of 
Technology\\
CH-8092 Zurich, Switzerland \\
mkalk@math.ethz.ch}



\date{} 
\pagestyle{empty}



\begin{document}
\maketitle

\centerline{\bf Abstract}

\medskip
The problem of converting parametrically defined varieties into 
their implicit form is of great importance in geometric modeling.
It essentially is an elimination problem and therefore 
techniques like resultants, Gr\"obner bases or characteristic sets
can be applied.
Gr\"obner bases have the important 
advantage that
they can be used for solving the implicitization problem in full 
generality. On the other hand elimination by means of Buchberger's 
Gr\"obner bases algorithm is often very time-consuming.

In this talk we present an alternative approach: 
implicitization by Gr\"ob\-ner basis conversion.
The idea underlying Gr\"obner basis
conversion is the following. It is well-known that the computing
time of a Gr\"obner basis computation 
heavily depends on the chosen order. Instead of computing 
a Gr\"obner basis $G_1$
with respect to a ``slow'' order directly
we could compute a Gr\"obner basis $G_2$
with respect to a ``fast'' order first and then convert $G_2$ to $G_1$.
If the conversion can be done efficiently this method will be faster than
the direct approach. 

We present three conversion algorithms. 
The FGLM-algorithm solves the conversion problem by linear algebra techniques.
In the Gr\"obner walk algorithm the conversion is done in several 
steps following a
path in the Gr\"obner fan of the given ideal. The third algorithm uses the
knowledge of the
Hilbert-Poincar\'e series of an ideal $I$ to speed up Gr\"obner bases 
computations for $I$.

Gr\"obner basis conversion is well applicable to the implicitization
problem since the parametric equations form ``almost'' a Gr\"obner basis.
We investigate how these three conversion 
algorithms can be applied.
\vspace{2 mm}
\newline
{\bf keywords:} rational parametric varieties, implicitization, Gr\"obner 
basis, basis conversion.
\end{document}
